% 图5: 端到端RDMA数据传输路径图
\begin{tikzpicture}[scale=0.85, transform shape,
    box/.style={rectangle, draw, rounded corners, minimum width=2cm, minimum height=0.8cm, align=center, font=\small},
    mem/.style={rectangle, draw, fill=blue!#1, minimum width=1.5cm, minimum height=0.6cm, font=\tiny},
    arrow/.style={->, >=stealth, thick},
    label/.style={font=\small\bfseries},
    stage/.style={rectangle, draw, fill=gray!#1, minimum width=0.6cm, minimum height=3cm},
    nocopy/.style={rectangle, draw=red, dashed, thick, inner sep=2pt}
]

% === 传统TCP/IP vs RDMA对比 ===

% 左侧标题
\node[label] at (-6, 6) {\large \textbf{传统TCP/IP路径}};
\node[font=\scriptsize, text=red] at (-6, 5.5) {多次拷贝 · 内核参与 · 高延迟};

% TCP/IP 数据路径
\node[mem={10}] (tcp-app) at (-9, 4) {应用缓冲区};
\node[mem={20}] (tcp-user) at (-9, 2.8) {用户空间};
\node[mem={30}] (tcp-kernel) at (-9, 1.6) {内核缓冲区};
\node[mem={40}] (tcp-nic) at (-9, 0.4) {网卡缓冲区};
\node[box, fill=gray!30] (tcp-net) at (-9, -0.8) {网络};

% 拷贝箭头
\draw[arrow, red!60] (tcp-app) -- (tcp-user) node[midway, right, font=\tiny] {拷贝1};
\draw[arrow, red!60] (tcp-user) -- (tcp-kernel) node[midway, right, font=\tiny] {拷贝2};
\draw[arrow, red!60] (tcp-kernel) -- (tcp-nic) node[midway, right, font=\tiny] {拷贝3};
\draw[arrow, red!60] (tcp-nic) -- (tcp-net);

% CPU参与标注
\node[rectangle, draw=orange, fill=orange!10, font=\tiny] at (-6.5, 3.4) {CPU参与};
\node[rectangle, draw=orange, fill=orange!10, font=\tiny] at (-6.5, 2.2) {CPU参与};
\node[rectangle, draw=orange, fill=orange!10, font=\tiny] at (-6.5, 1) {CPU参与};

% 延迟标注
\node[font=\scriptsize, text=red] at (-6.5, -1.5) {延迟: 10-100 μs};

% === 中间分隔线 ===
\draw[dashed, thick, gray] (0, -2) -- (0, 6.5);
\node[font=\bfseries, text=gray] at (0, 6.8) {vs};

% === 右侧: RDMA路径 ===
\node[label] at (6, 6) {\large \textbf{RDMA零拷贝路径}};
\node[font=\scriptsize, text=green!60!black] at (6, 5.5) {零拷贝 · 内核旁路 · 微秒级延迟};

% 发送方
\node[label, font=\small] at (2.5, 4.5) {发送方};
\node[mem={20}] (rdma-app-s) at (2.5, 3.5) {应用内存\\(Memory Region)};
\node[box, fill=green!20] (rdma-rnic-s) at (2.5, 2) {RNIC};
\node[rectangle, draw, fill=green!10, minimum width=1.5cm, minimum height=0.5cm, font=\tiny] (rdma-sq) at (2.5, 1.2) {发送队列(SQ)};
\node[rectangle, draw, fill=green!10, minimum width=1.5cm, minimum height=0.5cm, font=\tiny] (rdma-cq-s) at (2.5, 0.6) {完成队列(CQ)};

% 网络
\node[box, fill=gray!30, minimum width=2cm] (rdma-net) at (6, 2) {RoCEv2网络\\UDP/IP};

% 接收方
\node[label, font=\small] at (9.5, 4.5) {接收方};
\node[mem={20}] (rdma-app-r) at (9.5, 3.5) {应用内存\\(Memory Region)};
\node[box, fill=green!20] (rdma-rnic-r) at (9.5, 2) {RNIC};
\node[rectangle, draw, fill=green!10, minimum width=1.5cm, minimum height=0.5cm, font=\tiny] (rdma-rq) at (9.5, 1.2) {接收队列(RQ)};
\node[rectangle, draw, fill=green!10, minimum width=1.5cm, minimum height=0.5cm, font=\tiny] (rdma-cq-r) at (9.5, 0.6) {完成队列(CQ)};

% 绘制RDMA数据流
% 发送方内部
\draw[arrow, green!60!black] (rdma-app-s) -- (rdma-rnic-s) node[midway, right, font=\tiny] {直接DMA};
\draw[arrow, green!60!black] (rdma-rnic-s) -- (rdma-net) node[midway, above, font=\tiny] {数据包};
\draw[arrow, green!60!black] (rdma-net) -- (rdma-rnic-r) node[midway, above, font=\tiny] {数据包};
\draw[arrow, green!60!black] (rdma-rnic-r) -- (rdma-app-r) node[midway, right, font=\tiny] {直接DMA};

% 零拷贝标注
\node[nocopy, fit={(rdma-app-s) (rdma-rnic-s)}] {};
\node[nocopy, fit={(rdma-app-r) (rdma-rnic-r)}] {};
\node[font=\tiny, text=red] at (2.5, 4.1) {零拷贝区域};
\node[font=\tiny, text=red] at (9.5, 4.1) {零拷贝区域};

% 延迟标注
\node[font=\scriptsize, text=green!60!black] at (6, -0.5) {延迟: 1-2 μs (降低10倍+)};

% === 底部: 三种RDMA操作类型 ===
\node[label] at (0, -2.5) {\textbf{RDMA三种操作类型}};

% Send/Recv
\node[rectangle, draw, fill=blue!10, text width=3.5cm, align=center, font=\small] at (-4, -4) {
    \textbf{Send/Recv}\\
    传统消息传递\\
    接收方需预提交缓冲区\\
    双边操作
};

% RDMA Write
\node[rectangle, draw, fill=green!10, text width=3.5cm, align=center, font=\small] at (0, -4) {
    \textbf{RDMA Write}\\
    单边写操作\\
    无需接收方CPU参与\\
    发送方直接写入
};

% RDMA Read
\node[rectangle, draw, fill=orange!10, text width=3.5cm, align=center, font=\small] at (4, -4) {
    \textbf{RDMA Read}\\
    单边读操作\\
    无需接收方CPU参与\\
    发送方直接读取
};

% 底部总结
\node[rectangle, draw, fill=gray!10, text width=14cm, align=center, font=\small] at (0, -5.5) {
    \textbf{核心突破}: RDMA通过"内核旁路+零拷贝+CPU卸载"三管齐下，将网络延迟从TCP/IP的10-100微秒\\
    降低到1-2微秒，使GPU能够以内存速度进行跨节点通信。
};

\end{tikzpicture}
